\scanpage{63}

\begin{tabular}{|l}Cas des $\unsure{n|m}$-polygônes réguliers, \deleted{\unsure{de}} $n$ premier à la\\ caractéristique ($k$ alg clos)\end{tabular}

$S$ ensemble des sommets, $A$ ens des arêtes, $E$ espace affine ambiant.

Le barycentre $o$ de $S$ est fixé par le groupe $G$ des automorphismes de $\Pi$, de sorte que $E$ devient un vectoriel. Soit $G^+$ le sous-groupe des ``rotations'' de $\Pi$, $u$ un des $2$ générateurs privilégiés, correspondant à une orientation \inserted{combinatoire} de $\Pi$, \deleted{Considérons} \inserted{soit $s_0\in S$} \deleted{les deux valeurs propres de $u$, \unsure{puis}} [la donnée de $(s_0,u)$ équivaut à celle d'un repère combinatoire, i.e. d'un épinglage combinatoire] $\sigma$ l'autom $\ne\mathrm{id}$ tel que $\sigma s_0=s_0$, de sorte que

\[\left\{\begin{aligned}
&\sigma^2=\mathrm{id},\quad \sigma u\sigma^{\unsure{-1|\prime\prime}}=u^{-1}\\
&u^n=1\qquad\text{(pour mémoire)}
\end{aligned}\right.\]

Soient $\zeta,\zeta'$ les valeurs propres de $u$. Si on avait $\zeta=\zeta'$, on aurait $u=\zeta(\mathrm{id}_E+v)$, $v$ nilpotent, \inserted{i.e. $v^2=0$} et $u^n=\mathrm{id}\Longrightarrow\zeta^n=1$, $1+nv=0$ i.e. $nv=0$, d'où (comme $n1_{\unsure{k|K}}\ne0$ par hyp) $v=0$. donc $u=\zeta\,\mathrm{id}_E$, donc les $s_i=u^is_0$ seraient sur une $\overline{\mathrm m}$ droite, \uline{absurde}. Donc $\zeta\ne\zeta'$, soient $D_\zeta,D_{\zeta'}$ les droites propres. \deleted{\unsure{Mais}} Comme $u\sim u^{-1}$ l'un des

% archive sheet number: 62
